% Part: methods
% Chapter: induction
% Section: induction-on-N

\documentclass[../../../include/open-logic-section]{subfiles}

\begin{document}

\olfileid{mth}{ind}{inN}

\olsection{$\Nat$-এর উপর আরোহ}

সবচেয়ে সহজ রূপে আরোহ হলো সব স্বাভাবিক সংখ্যা সম্পর্কে
ফল প্রমাণের কৌশল। এর ভিত্তি হলো: $0$ থেকে শুরু করে
বারবার~$1$ যোগ করলে শেষ পর্যন্ত প্রত্যেক স্বাভাবিক সংখ্যায়
পৌঁছানো যায়। তাই প্রতিটি সংখ্যার জন্য কিছু সত্য প্রমাণ করতে
আমরা (1)~দেখাতে পারি যে তা $0$-এর জন্য সত্য, এবং
(2)~দেখাতে পারি যে কোনো সংখ্যা~$n$-এর জন্য সত্য হলে
পরের সংখ্যা~$n+1$-এর জন্যও সত্য। ``সংখ্যা~$n$-এর ধর্ম
$P$ আছে'' কথাটি $P(n)$ দিয়ে সংক্ষেপ করি (একইভাবে
``সংখ্যা~$k$-এর ধর্ম $P$ আছে'' কথাটি $P(k)$ দিয়ে)।
তাহলে সব $n \in \Nat$-এর জন্য $P(n)$ প্রমাণের একটি
আরোহ-প্রমাণে থাকবে:
\begin{enumerate}
\item $P(0)$-এর প্রমাণ, এবং
\item যেকোনো~$k$-এর জন্য $P(k)$ হলে $P(k+1)$ হয়—এর প্রমাণ।
\end{enumerate}
ব্যাপারটি স্পষ্ট করতে ধরা যাক, (1) এবং~(2) উভয়ই আমাদের
আছে। (1) থেকে $P(0)$ সত্য। (2) থেকে বিশেষভাবে জানি,
$P(0)$ হলে $P(0+1)$, অর্থাৎ $P(1)$। সাধারণ বক্তব্য
``যেকোনো~$k$-এর জন্য $P(k)$ হলে $P(k+1)$''-এ~$k$-এর
জায়গায়~$0$ বসিয়ে এটি পাই। অতএব মোডাস পোনেন্সে~$P(1)$।
% BN-SRC-483: the induction step quantifies k, so substitute for k.
(2)-এ এবার $k$-এর জায়গায় $1$ নিলে পাই: $P(1)$ হলে~$P(2)$।
$P(1)$ ইতিমধ্যে প্রতিষ্ঠিত, তাই মোডাস পোনেন্সে~$P(2)$।
এভাবেই চলতে থাকে। যেকোনো সংখ্যা~$n$-এর জন্য এই ধাপ
$n$~বার প্রয়োগ করে শেষ পর্যন্ত~$P(n)$-এ পৌঁছাই। ফলে
(1) ও~(2) একসঙ্গে যেকোনো $n \in \Nat$-এর জন্য $P(n)$
প্রতিষ্ঠা করে।

একটি উদাহরণ দেখা যাক। $n$টি পাশা ছুড়ে কত রকম মোট
পাওয়া যায়, তা জানতে চাই। অদ্ভুত মনে হলেও $0$টি পাশা
দিয়ে শুরু করি। কোনো পাশাই না থাকলে একটিই সম্ভাব্য
``ছোড়ার'' ফল: একটিও ফোঁটা নেই, মোট~$0$। সুতরাং ভিন্ন
সম্ভাব্য মোটের সংখ্যা~$1$। একটি পাশা, অর্থাৎ $n=1$ হলে
$1$ থেকে~$6$ পর্যন্ত ছয়টি সম্ভাব্য মান। দুটি পাশায় $2$
থেকে $12$ পর্যন্ত যেকোনো মোট পাওয়া যায়, অর্থাৎ $11$টি
সম্ভাবনা। তিনটি পাশায় $3$ থেকে $18$ পর্যন্ত যেকোনো সংখ্যা
পাওয়া যায়, অর্থাৎ $16$টি সম্ভাবনা। $1$, $6$, $11$, $16$:
একটি ধারা দেখা যাচ্ছে—উত্তর হয়তো $5n+1$? ভিন্ন সম্ভাব্য
মোটের সংখ্যা অবশ্যই $5n+1$-এর বেশি হতে পারে না। কারণ $n$টি পাশায়
সর্বনিম্ন মোট $n$ (সবগুলিতে $1$) এবং সর্বোচ্চ $6n$
(সবগুলিতে $6$); এই দুই সীমার মধ্যে মোট $5n+1$টি সংখ্যা।

\begin{thm}
  $n$টি পাশা দিয়ে $n$ থেকে $6n$ পর্যন্ত সম্ভাব্য $5n+1$টি
  মানের সবগুলিই পাওয়া যায়।
\end{thm}

\begin{proof}
$P(n)$ দিয়ে বোঝাই: ``$n$টি পাশা দিয়ে $n$ থেকে~$6n$
পর্যন্ত যেকোনো সংখ্যা পাওয়া সম্ভব।'' আরোহ প্রয়োগ করতে
প্রমাণ করি:
% BN-SRC-838 TARGET-CORRECTION-BEGIN
\par\noindent\textbf{পাঠ-সংশোধন।} এখানে ভিন্ন মোটের সংখ্যার ঊর্ধ্বসীমা বলা হচ্ছে; উল্টো
অসমতা নয়। শূন্য পাশার একমাত্র মোট শূন্য—আগের
অনুচ্ছেদেই এই ক্ষেত্র যাচাই হয়েছে। তাই পরের $P(1)$-ভিত্তি
ও উত্তরসূরি-ধাপ ধনাত্মক সংখ্যাগুলি প্রমাণ করলে, আগে
যাচাই করা $P(0)$-সহ সব স্বাভাবিক সংখ্যার দাবি পাওয়া যায়।
ধনাত্মক অংশে ধাপের $k$ অন্তত এক; শূন্যকে বাদ দিয়ে
উপপাদ্যের পরিধি সংকুচিত করা হয়নি।
% BN-SRC-838 TARGET-CORRECTION-END
\begin{enumerate}
\item \emph{আরোহের ভিত্তি} $P(1)$: একটি পাশা দিয়ে $1$
  থেকে $6$ পর্যন্ত যেকোনো সংখ্যা পাওয়া যায়।
\item \emph{আরোহের ধাপ}: প্রত্যেক $k$-এর জন্য $P(k)$ হলে~$P(k+1)$।
\end{enumerate}

(1) প্রমাণে $6$-পিঠের একটি পাশা দেখলেই হয়। তার ছয়টি
পিঠে $1$ থেকে~$6$ পর্যন্ত প্রতিটি সংখ্যা একবার করে আছে।
তাই একটি পাশা দিয়ে $1$ থেকে~$6$ পর্যন্ত যেকোনো সংখ্যা
পাওয়া সম্ভব।

(2) প্রমাণ করতে শর্তাধীন বক্তব্যটির পূর্বপক্ষ, অর্থাৎ
$P(k)$, ধরে নিই। এই অনুমানকে \emph{আরোহের অনুমান}
বলা হয়। এটি কাজে লাগিয়ে~$P(k+1)$ প্রমাণ করি। কঠিন
বিষয় হলো, $k+1$টি পাশা ছুড়ে পাওয়া মানগুলিকে $k$টি
পাশায় পাওয়া মান এবং বাড়তি $k+1$-তম পাশার মানের
সাহায্যে ভাবা। আরোহের অনুমান ব্যবহার করতে চাইলে ঠিক
এ কাজটিই করতে হবে।

আরোহের অনুমান বলে, $k$টি পাশায় $k$ থেকে~$6k$ পর্যন্ত
যেকোনো সংখ্যা পাওয়া যায়। $(k+1)$-তম পাশায়~$1$
পড়লে মোটে $1$ যোগ হয়। তাই $k$টি পাশার সঙ্গে বাড়তি
$(k+1)$-তম পাশায়~$1$ পেলে $k+1$ থেকে $6k+1$ পর্যন্ত যেকোনো মান
পাওয়া যায়। বাকি থাকে $6k+2$ থেকে $6k+6$ পর্যন্ত
মানগুলি। প্রথম $k$টি পাশার প্রত্যেকটিতে $6$ এবং
$(k+1)$-তম পাশায় $2$ থেকে $6$-এর মধ্যে একটি সংখ্যা
পেলে সেগুলিও পাওয়া যায়। সব মিলিয়ে $k+1$টি পাশায়
$k+1$ থেকে $6(k+1)$ পর্যন্ত প্রতিটি সংখ্যা পাওয়া যায়।
অর্থাৎ আরোহের অনুমান~$P(k)$ ব্যবহার করে~$P(k+1)$
প্রমাণ করেছি।
\end{proof}

অনেক সময় যে বস্তুগুলির অনুক্রম (সংখ্যা, সেট ইত্যাদি)
সম্পর্কে কিছু প্রমাণ করতে চাই, সেই অনুক্রমই ``আরোহীভাবে''
সংজ্ঞায়িত: $n$-তম বস্তুর সাহায্যে $(n+1)$-তম বস্তুর
সংজ্ঞা দেওয়া হয়। তখন প্রায়ই আরোহ ব্যবহার করি। যেমন,
~$n$ পর্যন্ত স্বাভাবিক সংখ্যাগুলির যোগফল~$s_n$ সংজ্ঞায়িত
করতে পারি এভাবে:
\begin{align*}
  s_0 & = 0\\
  s_{n+1} & = s_n + (n+1)
\end{align*}
এই সংজ্ঞা থেকে পাই:
\begin{align*}
  s_0 & = 0,\\
  s_1 & = s_0 + 1 && = 1,\\
  s_2 & = s_1 + 2 && = 1 + 2 = 3\\
  s_3 & = s_2 + 3 && = 1 + 2 + 3 = 6, \text{ ইত্যাদি।}
\end{align*}
এখন আরোহ দিয়ে প্রমাণ করতে পারি যে $s_n = n(n+1)/2$।

\begin{prop}
  $s_n = n(n+1)/2$।
\end{prop}

\begin{proof}
  দেখাতে হবে (1) $s_0 = 0\cdot(0 + 1)/2$ এবং (2) $s_k =
  k(k+1)/2$ হলে $s_{k+1} = (k+1)(k+2)/2$। (1) স্পষ্ট।
  (2) প্রমাণে আরোহের অনুমান নিই: $s_k = k(k+1)/2$।
  এটি ব্যবহার করে দেখাতে হবে $s_{k+1} = (k+1)(k+2)/2$।

  $s_{k+1}$ কত? সংজ্ঞা অনুসারে $s_{k+1} = s_k + (k+1)$।
  আরোহের অনুমানে $s_k = k(k+1)/2$। আগের সমীকরণে তা
  বসিয়ে ভগ্নাংশের সামান্য হিসাব করলেই হয়:
  \begin{align*}
    s_{k+1} & = \frac{k(k+1)}{2} + (k+1) = {}\\
    & = \frac{k(k+1)}{2} + \frac{2(k+1)}{2} = {}\\
    & = \frac{k(k+1) + 2(k+1)}{2} = {}\\
    & = \frac{(k+2)(k+1)}{2}.
  \end{align*}
\end{proof}

এখানে গুরুত্বপূর্ণ শিক্ষা হলো: কোনো আরোহীভাবে সংজ্ঞায়িত
অনুক্রম $a_n$ সম্পর্কে কিছু প্রমাণ করতে হলে আরোহই
স্বাভাবিক পথ। অনুক্রমটি এভাবে সংজ্ঞায়িত না হলেও
(যেমন পাশার মোটের সম্ভাবনার ক্ষেত্রে), কোনোভাবে~$k+1$
ক্ষেত্রকে~$k$ ক্ষেত্রের সঙ্গে যুক্ত করতে পারলে আরোহ
ব্যবহার করা যায়।

\end{document}
